\section{Abstract}
Preference-conditioned multi-objective reinforcement learning (MORL) aims to
represent several objective trade-offs with one policy. This project implements
a vector Q-network and investigates whether difficulty, ensemble disagreement,
coverage/recency, and observed learning progress can guide training allocation.
The recovered primary experiment comprises 30 runs: two environments, three
allocation arms, and five paired training seeds, with 500,000 interactions per
run. All 300 retained checkpoint-evaluation records were numerically re-audited,
covering 28,800 validation episodes. Minecart produced zero ore in both ore
objectives throughout its 14,400 primary validation episodes. The comparative
qualification therefore failed; adaptive superiority is not established.
We present the implementation, descriptive results, failure analysis, recovery
limitations, and a real inference demonstration. A separate predeclared Envelope
backbone feasibility screen is reported independently. A later uniform on-policy
Minecart confirmation passed both task profiles in both inference modes across
three fresh training seeds. Lunar competence and adaptive-allocation benefit
remain unestablished. The original full baseline/ablation matrix and final-test
study remain incomplete.

\section{Problem and research questions}
\subsection{Motivation and scope}
In ordinary reinforcement learning, an agent maximizes one scalar reward.
MORL instead provides $r_t\in\mathbb{R}^d$, with potentially conflicting
objectives. A user specifies $w\in\Delta^{d-1}$, where $w_i\geq0$ and
$\sum_iw_i=1$. The discounted return is
\[
 G=\sum_{t=0}^{T-1}\gamma^t r_t,\qquad U_w(G)=w^\top G.
\]
A preference-conditioned policy receives both observation and preference,
allowing users to change trade-offs without retraining a separate model.
This does not guarantee that every requested trade-off is learned.

Training time is limited. A difficult region, an uncertain region, and a region
with improving performance are different reasons to allocate additional effort.
The project asks which signals predict later improvement (RQ1), whether adaptive
allocation improves learning over uniform and matched-random controls (RQ2),
whether effects arise at rollout or update time (RQ3), and whether results
transfer across preferences or dynamics (RQ4). The executed evidence does not
fully answer all four questions. The undergraduate deliverable is an implemented,
auditable investigation and demonstration, with unresolved questions stated
explicitly; acceptance against the lecturer's rubric remains external to this report.

\subsection{Environments}
\texttt{mo-lunar-lander-v3}\footnote{\url{https://mo-gymnasium.farama.org/environments/mo-lunar-lander/}}
supplies four objectives covering landing outcome,
shaping, and two engine costs. \texttt{minecart-v0}\footnote{\url{https://mo-gymnasium.farama.org/environments/minecart/}}
supplies two delivered-ore
objectives and a fuel-cost objective. Minecart requires a multi-step sequence
of navigation, mining and delivery; a low fuel cost alone does not establish
ore collection. Both environments have discrete actions and vector rewards.
Frozen course protocols use $\gamma=0.99$ and pilot-derived static reward scales
and hypervolume reference points. They are not recalibrated using final-test data.

\section{Research background}
Dynamic-weight Q-learning~\cite{abels2019dynamic-weights} and Envelope
Q-learning~\cite{yang2019envelope} are direct antecedents of conditioned
multi-objective value learning. PCN conditions a policy on desired return and
horizon~\cite{reymond2022pcn}; it differs from this project's preference-conditioned
Q-network. Preference conditioning is therefore not a novelty claim.
PG-MORL~\cite{xu2020pgmorl} and GPI prioritization~\cite{alegre2023gpi-prioritization}
already direct effort toward promising trade-offs. The project's narrower
interest is a controlled comparison of allocation signals and their failure modes.

Value disagreement motivates exploration~\cite{zhang2020value-disagreement},
while ALP-GMM~\cite{portelas2020alp-gmm} and CURIOUS~\cite{colas2019curious}
motivate competence-progress curricula. Absolute change can include forgetting;
it is not automatically signed improvement. A shared representation also means
our head disagreement is not independently calibrated epistemic uncertainty.
These sources motivate mechanisms, not claims that they work in this implementation.

\section{Implemented system}
\begin{figure}[htbp]
\centering
\resizebox{\textwidth}{!}{\begin{tikzpicture}[node distance=10mm and 9mm]
\node[block] (pref) {Preference sampler\\$w\sim p_t(w)$};
\node[block,right=of pref] (env) {Rollout + environment\\vector reward $r$};
\node[block,right=of env] (replay) {Replay buffer\\transitions + provenance};
\node[block,below=18mm of replay] (q) {Masked-head PCQN\\vector Double-DQN update};
\node[block,left=of q] (scores) {Curriculum diagnostics\\gap / disagreement\\coverage / progress};
\node[block,left=of scores] (alloc) {Budgeted allocation\\uniform-floor mixture};
\draw[flow] (pref)--(env); \draw[flow] (env)--(replay);
\draw[flow] (replay)--(q); \draw[flow] (q)--(scores);
\draw[flow] (scores)--(alloc); \draw[flow] (alloc)--(pref);
\node[below=11mm of scores,align=center,font=\sffamily\small,text=Muted]
{Separate validation selects checkpoints; final-test remains locked.\\
Controls share the preference pool and diagnostic budget.};
\end{tikzpicture}
}
\caption{Implemented allocation loop. Validation is separate from training;
checkpoint selection and final test remain subject to scientific gates.}
\end{figure}

\subsection{Vector value function and update}
The network receives $(s,w)$ and emits a vector value for each action, using
three masked heads over a shared trunk. Greedy inference uses
\[
 a^*(s,w)=\arg\max_a w^\top\frac1K\sum_{k=1}^K Q_{\theta,k}(s,a,w).
\]
Each head has its own Double-DQN action selection in the Bellman target:
\[
 a'_k=\arg\max_a w^\top Q_{\theta,k}(s',a,w),\quad
 y_k=r+\gamma(1-\mathrm{terminated})Q_{\bar\theta,k}(s',a'_k,w).
\]
Masked vector Huber losses update the online network; a target network supplies
continuation values. True termination disables bootstrapping, whereas time-limit
truncation ends collection without automatically suppressing continuation.
Replay retains the vector reward, behavior preference, both end flags, and a
bootstrap mask. Analytic unit tests cover these contracts.

\subsection{Allocation mechanism and controls}
An archive-relative conditioning gap estimates current utility shortfall.
Masked-head disagreement measures variation across head predictions.
Coverage/recency penalizes neglected candidate regions. Absolute utility change
between diagnostic measurements estimates observed learning progress.
Normalized scores form a softmax distribution mixed with a uniform floor:
\[
 p_t(w_i)=(1-\eta)\operatorname{softmax}(z_t/\tau)_i+\eta/N.
\]
The uniform floor preserves exploration. The implementation distinguishes
rollout-only, update-only, and joint allocation, but the recovered primary block
only compares uniform, matched-random and full joint allocation. The remaining
factorial ablations must not be implied by the existence of their code paths.

Uniform uses the same fixed candidate pool. Matched-random retains the
diagnostic schedule but replaces adaptive scores with random scores. Interaction
and gradient-update budgets are matched; actual evaluation environment steps
and wall-clock can differ because episode lengths differ. There is no claim
that every source of compute cost is exactly equal.

\section{Experimental design and evidence provenance}
\subsection{Recovered primary block}
Each of three arms was run on both environments with five paired seeds
(20262001--20262005): 30 jobs and 15 million training interactions in total.
Each job used 500,000 interactions and 123,751 gradient updates, with ten saved
checkpoints at 50,000-step intervals. Nominal validation used 32 fixed preferences
and three episodes per preference. The aggregate retains all 300 checkpoint
records, 28,800 episode vectors and 18,304,760 evaluation steps.

The latter count covers the retained checkpoint re-evaluations, not the entire
training-time evaluation cost. Historical final counters separately record
36,758,477 validation steps and 19,531,144 curriculum-evaluation steps.
Training, updates, curriculum and validation are separate accounting categories.
Repeated episodes across checkpoints are not independent training replicates;
the training seed is the experimental unit.

\subsection{Metrics and data units}
Expected utility is the mean of $w_i^\top\bar G_i$ across fixed preferences.
Hypervolume measures the union of regions dominated by return vectors above a
fixed maximization reference point. Both metrics use discounted protocol-scaled
returns. A nonzero hypervolume can occur without completing Minecart's ore task.
The field named \texttt{raw\_episode\_returns} is undiscounted
\emph{protocol-scaled} reward, not native environment units.

\subsection{Recovery limits}
The latest recovered pushed research snapshot is \texttt{e223362}, dated
18 October 2026. The user recalls local code edits and experiments through
22 October on the subsequently lost laptop; those changes have not been recovered.
The aggregate contains numerical episode evidence, but all 300 referenced
historical checkpoint files are missing. The 1 November re-audit verifies
retained tensors, preferences, metadata, protocol/source bindings and metrics.
It is not a new rollout reproduction or a new verification of missing model files.

\section{Results and failure analysis}
\subsection{Fixed-budget descriptive results}
Table~\ref{tab:primary} uses the fixed 500,000-step checkpoint and reports
mean $\pm$ sample standard deviation across five seeds. These are validation
diagnostics, not selected final-test results or qualified superiority estimates.
\begin{table}[htbp]
\centering\small
\caption{Recovered primary validation at 500,000 interactions; mean $\pm$ sample SD, five seeds.}
\label{tab:primary}
\begin{tabular}{llrr}
\toprule
Environment & Arm & Hypervolume & Expected utility\\
\midrule
LunarLander & Uniform & $4.0609\pm0.3293$ & $-0.0758\pm0.0311$\\
LunarLander & Matched-random & $4.6750\pm0.4427$ & $-0.0227\pm0.0684$\\
LunarLander & Full adaptive & $4.1764\pm0.3228$ & $-0.0458\pm0.0423$\\
Minecart & Uniform & $0.000253\pm0.000565$ & $-0.6912\pm0.1455$\\
Minecart & Matched-random & $0.000561\pm0.000768$ & $-0.5857\pm0.1577$\\
Minecart & Full adaptive & $0.000732\pm0.000671$ & $-0.6302\pm0.0443$\\
\bottomrule\end{tabular}\end{table}


\begin{figure}[H]
\centering
\includegraphics[width=\textwidth]{../../outputs/course-report-latex/recovered_learning_curves.pdf}
\caption{All retained primary checkpoints: mean and sample standard deviation
across five training seeds. Shading is seed variation, not a confidence interval.
These are validation diagnostics; no best-checkpoint selection is applied.}
\end{figure}

The full adaptive arm is not numerically best on LunarLander in this summary.
Neither this observation nor a small Minecart metric difference establishes
superiority, equivalence, or a general mechanism. All 14,400 Minecart primary
validation episodes have zero positive return in both ore objectives. The
study was marked \texttt{failed\_degenerate}; its checkpoint-selection,
remaining matrix-v1 runs and final-test workflow remain blocked.

\subsection{Why more steps alone are not a completion rule}
Separate remediation studies investigated the backbone before attempting an
allocation conclusion. PER-only v2 ran three one-million-interaction seeds;
its all-seeds qualification failed. F1 tested successful-episode replay and
achieved-return relabeling in six paired jobs. Its source audit covered 60
checkpoints, with 9,600 validation episodes recorded historically. The A1R1
variant had zero positive ore episodes at every checkpoint (4,800 episodes).
The A0R0 comparator recorded 79 positive cases per ore objective across 4,800
repeated validation episodes, not 79 independent training runs. F1 failed its
retention screen. These separate studies are not merged into Table~\ref{tab:primary}.
Unlike the primary aggregate, their complete raw tensors and models were not
recovered locally; interpretation relies on retained audit records.

Sparse delivery, temporal credit assignment and preference alignment remain
candidate explanations. The evidence does not identify a unique cause or show
that all replay methods fail. Repeatedly extending a failed experiment until
a positive result appears would change the question after seeing outcomes.

\subsection{Prospective Envelope backbone feasibility}
The predeclared screen used MORL-Baselines 1.3.0 Envelope and the pinned
upstream Minecart example, with the repository's frozen $\gamma=0.99$ and reward
scales. Each of three seeds ran exactly 100,000 training interactions and
99,901 gradient updates. Gaussian behavior/update preferences were retained;
this is not a fixed-pool or compute-matched allocation comparator. Validation
used all 32 course preferences, three episodes each, at fixed steps 50,000 and
100,000. Training RNGs and environments were isolated from validation. The
original CUDA smoke exposed a replay-index transport bug; a tested host-index
bridge preserved the sampled indices/priorities and allowed the pinned learner
to run. No learning rule or hyperparameter was tuned on these validation results.

\begin{table}[htbp]
\centering\small
\caption{Envelope feasibility: positive ore episodes out of 96 fixed cases at each checkpoint.}
\begin{tabular}{rrrrrr}
\toprule
Seed & Steps & Ore 1 & Ore 2 & HV & Expected utility\\
\midrule
20261121 & 50,000 & 0 & 0 & 0.001477 & -0.5444\\
20261121 & 100,000 & 0 & 0 & 0.001448 & -0.5512\\
20261122 & 50,000 & 0 & 0 & 0.001421 & -0.5868\\
20261122 & 100,000 & 2 & 2 & 0.001359 & -0.7790\\
20261123 & 50,000 & 2 & 2 & 0.001283 & -0.6058\\
20261123 & 100,000 & 0 & 0 & 0.001420 & -0.5967\\
\bottomrule\end{tabular}\end{table}

The completed cohort used 300,000 training interactions, 299,703 updates and 530,664 separate validation environment steps across 576 episodes. Final-test steps were zero. At least one final checkpoint expressed positive ore in some validation cases. The per-seed counts above delimit that observation; they do not establish robust task coverage or an adaptive-allocation benefit. All model files and validation tensors passed hash, shape, finite-value, budget and metric checks. This numerical audit is not an independent semantic review. Neither outcome unlocks the failed matrix-v1/F1 or final-test workflow.


\subsection{Task-level positive control and protocol review}
The presentation default is the validation preference nearest the uniform
simplex point, not a checkpoint or seed chosen for successful behavior.
Equal preference coordinates multiply scaled rewards; they do not guarantee
equal task importance. A separate audit of eight preferences and five environment
seeds per environment found zero ore collection/delivery in all 40 Minecart
cases. Slider sensitivity therefore demonstrates input conditioning, while
task success requires a separate check.

We added a positive control on ten new paired Lunar development seeds:
the frozen model at default and landing-heavy weights, and Gymnasium's
unmodified handwritten heuristic. This is a fixed controller, not learned MORL.
Native per-step reward, state and action traces are retained.
\begin{table}[H]
\centering\small
\caption{Lunar positive control: ten paired environment seeds for one frozen
model, not independent training replicates. Pad x is an additional spatial diagnostic.}
\begin{tabular}{lrrr}
\toprule
Controller & Native +100 & +100 and pad x & Mean steps\\
\midrule
Model, default & 0/10 & 0/10 & 1000.0\\
Model, landing-heavy & 4/10 & 2/10 & 683.2\\
Handwritten heuristic & 8/10 & 7/10 & 216.9\\
\bottomrule
\end{tabular}
\end{table}
At the unchanged protocol $\gamma=.99$, the heuristic has greater
default-weight utility in 9/10 pairs (means .2340 versus .0574).
At landing-heavy weights it also exceeds the corresponding model in 9/10
pairs (means .0585 versus $-.1961$). Thus objective/task mismatch alone does
not explain the failure: a better controller exists under the same utility.
The native +100 condition checks the sleeping lander body, without an explicit
pad-x test; reporting only +100 would omit this distinction.

The full-text/source review also corrects budget comparability.
Abels et al.~\cite{abels2019dynamic-weights} report Minecart to one million
decision steps with a different optimizer and discount factor. The 100k
Envelope setup is a library example; Minecart is not an experimental task in
the original Envelope paper~\cite{yang2019envelope}.
The author PCN Minecart implementation~\cite{reymond2022pcn} instead uses a
deterministic environment and ten million interactions.
A later stochastic-Minecart architecture study uses four million interactions
and five seeds~\cite{terekhov2024architectures}. These are implementation
differences to disclose, not evidence that longer training alone fixes our learner.

A new development-only screen was frozen before execution: $n=1$ versus
$n=20$ vector returns, both environments, three paired training seeds, and
500k interactions per job (12 jobs, six million planned interactions).
All other learner settings, uniform preference allocation, static scales and
$\gamma=.99$ are matched. The $n=20$ choice is an unvalidated long-credit
hypothesis, not a paper-established Minecart default. Uncorrected off-policy
multi-step targets may introduce bias. Short repeated CUDA learning traces
matched exactly before main execution. Main training is currently ongoing.
No outcome or completed six-million-step budget is claimed in this report update.

Only final checkpoints can pass the predeclared development task screen:
each training seed needs at least 8/10 native +100-and-pad Lunar cases at
landing-heavy weights, or 5/10 ore deliveries at the declared Minecart
ore-seeking weights. Intermediate 100k/250k checkpoints describe learning and
cannot select a best model. The original allocation matrix and final-test
gates remain separate. Evidence and plans are indexed in
\path{docs/research/BACKBONE_RESEARCH_REVIEW_20261101.md}.


\subsection{From failed value policies to an on-policy backbone}
The completed n-step and subsequent relabel/loss screens failed their task gates;
the latter produced zero deliveries in 240 final task cases despite ore-positive
training episodes. These failures motivated a separate categorical actor with
a vector state-value critic, rather than extending a failed experiment.
The candidate uses a multibody architecture, clipped PPO, critic PopArt and
MDMM entropy control motivated by MOPPO~\cite{terekhov2024architectures}.
Unlike its objective-wise normalized actor advantage, this project scalarizes
the original fixed-scaled vector advantage before batch standardization, preserving
the protocol utility weights. It is a new candidate, not an exact reproduction.

A one-preference specialist and then a fresh uniform-32-preference conditional
pilot each ran 4 million interactions. Their final task gates passed. A separate
confirmation froze three fresh training seeds (20264311--20264313), new task
cases, 32 new Dirichlet preferences, and seven anchors reported separately.
Each seed used 4 million interactions, 250,080 actor and 500,160 critic updates.
Training allocation was uniform per episode; varying episode lengths imply
unequal interaction exposure per weight. No final-test data were accessed.

\begin{table}[htbp]
\centering\small
\caption{Final Minecart confirmation: deliveries on 10 new cases per profile.
The three training seeds are the replication units.}
\begin{tabular}{lcccc}\toprule
Seed & Sampled balanced & Sampled ore1 & Argmax balanced & Argmax ore1\\\midrule
20264311 &10/10&10/10&10/10&10/10\\
20264312 &10/10&10/10&10/10&10/10\\
20264313 &10/10&10/10&10/10&10/10\\\bottomrule
\end{tabular}
\end{table}
All final utilities also exceeded the frozen idle reference. Across training
seeds, expected utility on the 32 new development preferences was
$0.19835\pm0.03540$ for sampled inference and $0.17680\pm0.07376$ for argmax
(sample standard deviation, $n=3$). These averages are not a guarantee at every
preference. Anchors are excluded from these means.
\begin{figure}[htbp]\centering
\includegraphics[width=\textwidth]{../../outputs/course-report-latex/minecart_confirmation.pdf}
\caption{Confirmation learning trajectories. Each line is one training seed;
intermediate checkpoints are descriptive and final selection was fixed at 4M.}
\end{figure}

The source, checkpoint and numeric audit checked 3,288 development episodes,
372,775 evaluation steps and 480 task traces, including native rewards,
discounting, utility and exact preference exposure. This executor audit is not
an independent semantic-review PASS. The evidence supports the two-profile task
gate across three seeds under nominal geometry. It does not establish adaptive
allocation benefit, optimality or dynamics robustness. Earlier failed studies
and their budgets are retained; the PPO run is not a matched-budget comparison
against historical PCQN. The Lunar pilot and separate two-seed confirmation are reported
in the discussion; it is not part of this three-seed Minecart confirmation.

\evidence{research/evidence/minecart_onpolicy_confirmation_v1.json}


\section{Demonstration and reproducibility}
\subsection{Real inference demonstration}
The local application uses actual environment transitions and frozen network
inference. Reference and user panels begin with the same seed and checkpoint;
the slider changes preference input, not model weights. Pause/step, live
preference changes, original environment rendering and trace export expose
observations, actions and rewards for inspection. Live changes retain prior
trajectory history and therefore are not a fresh paired episode comparison.

The updated Minecart demo uses the final4M categorical PPO checkpoint from
confirmation seed20264311, selected by the fixed lowest-seed rule rather than
observed score. Inference uses actor argmax; the two panels share the same
weights and initial environment seed. The other two confirmation models and
all confirmation evidence are retained. The Lunar panel still uses the earlier
100k-interaction Uniform-PCQN replacement and is explicitly labelled as not
qualified for reliable landing. Neither panel is a recovered historical model.
The demonstration does not establish arbitrary preference or allocation benefit.

\subsection{Execution and quality assurance}
On Windows, run \path{START_COURSE_DEMO.bat} from the repository root, then
open \url{http://127.0.0.1:8765}. The package installer, Python runtimes,
downloads, caches, checkpoints and compiler outputs are stored on D:.
The Python 3.11 CPU environment is retained separately from the CUDA environment.
Baseline qualification is specific to its recorded runtime and dependency set;
it is not transferred to a different installation. Device execution does not
imply algorithmic success.

Historical verification recorded backend tests, 80 frontend assertions and an
extracted-package installation check. The 3 November PPO package additionally
passed 20 real presentation cases, portable/source actor equivalence, and Edge
desktop/mobile controls, rendering and PDF-link checks when launched outside
repository imports. Installation was not repeated for every report revision.
The Envelope runner separately tests validation RNG isolation and return/end
flags. These are engineering checks, not support for the research hypotheses.

\section{Discussion, limitations and conclusions}
\subsection{Answers supported by the executed evidence}
\begin{table}[htbp]\centering\small
\begin{tabularx}{\textwidth}{lY}\toprule
Question & Evidence-bounded answer\\\midrule
RQ1: predictive signals & Implemented; no complete prospective signed-improvement comparison.\\
RQ2: allocation benefit & Primary block failed task qualification; no supported adaptive advantage.\\
RQ3: rollout versus update & Code paths exist; the full factorial comparison remains unexecuted.\\
RQ4: generalization & Three-seed nominal Minecart task gate passes on new cases; limited development-weight coverage is described. Dynamics robustness is untested.\\\bottomrule
\end{tabularx}
\caption{Research-question coverage. Backbone confirmation does not answer all allocation questions.}
\end{table}

\subsection{Lunar completion rule}
A separate Lunar specialist pilot is frozen at four million interactions with
landing-heavy weights $(.7,.1,.1,.1)$. It uses categorical PPO with vector
GAE($\lambda=.95$), eight streams and 128-step fragments. Unlike Minecart's
six-action setting, its entropy schedule starts at $\log4$. Final-only task
assessment requires native +100 and lander x inside the pad on at least 8/10
new development cases in each inference mode. This single-seed pilot is still
running at the report build; no completed Lunar PPO outcome is included.
If it fails, the failure is reported without automatic extension. If it passes,
independent confirmation needs a separate protocol. Its planned budget is not
recorded as completed expenditure or multi-seed evidence.

\subsection{Limits of the current conclusion}
The deliverable establishes a working preference-conditioned value-learning
framework, controlled primary validation, an auditable negative Minecart outcome,
and a real inference demo. It does not establish that adaptive allocation is
better than uniform or random allocation. RQ1 lacks a complete four-signal
prospective signed-improvement study; RQ3 lacks the complete executed allocation
target ablations; RQ4 has no qualified dynamics-generalization result. Preference
generalization and dynamics generalization are separate claims.

Lost historical weights restrict independent rollout reproduction. The new
Envelope screen differs in learner, interaction budget and upstream preference
sampling; it cannot serve as a compute-matched comparison against the recovered
PCQN table. Its incomplete global seeding is an additional reproducibility limit.
The completed n-step and relabel/loss studies failed their task screens; the later three-seed PPO confirmation supports only its bounded Minecart task gate. Validation means and standard deviations do not replace a held-out
final test or an adequately powered study. A completed numerical review of an
earlier report returned WARN; the later revised-report reviewer hit its usage
limit. No final independent semantic PASS is claimed.

The defensible conclusion is that the backbone must demonstrate the task before
allocation benefits can be interpreted. This report documents a scoped course
investigation comprising implementation, primary validation, failure analysis
and bounded backbone screens, culminating in a three-seed Minecart task confirmation. The real inference demo exposes
preference-conditioned action selection. The original full proposal remains
partially fulfilled; a broader final
comparative study requires a separately frozen protocol and qualification,
not silent reopening of failed studies. The lecturer's final deadline and
minimum rubric have not yet been confirmed.

\appendix
\section{Evidence index and reproduction commands}
\begin{itemize}
\raggedright
\item \evidence{research/evidence/course_primary_validation_v1.json}: all retained primary episode vectors.
\item \evidence{research/evidence/course_primary_execution_audit_v1.json}: historical counters and source identities.
\item \evidence{configs/experiments/course_paired_matrix_v1.json}: original frozen design.
\item \evidence{research/qualification/minecart_pcqn_replay_f1.json}: failed F1 decision.
\item \evidence{configs/experiments/envelope_feasibility_v1.json}: new bounded feasibility plan.
\item \evidence{scripts/reaudit_recovered_validation.py}: numerical re-audit without environment rollouts.
\item \evidence{scripts/run_envelope_feasibility.py}: new validation-isolated Envelope runner.
\item \evidence{research/evidence/minecart_onpolicy_confirmation_v1.json}: completed three-seed Minecart task confirmation.
\item \evidence{docs/research/ONPOLICY_CONFIRMATION_RESULTS_20261103.md}: exact scope, budgets and limits.
\item \evidence{configs/experiments/lunar_onpolicy_specialist_v1.json}: running Lunar pilot, not a completed result.
\item \evidence{paper/REL301m/FINAL_REPORT_OUTLINE.md}: finalization outline and remaining evidence gates.
\item \evidence{paper/REL301m/CLAIM_EVIDENCE_MATRIX.md}: statements linked to retained data.
\item \evidence{docs/research/COURSE_FINAL_DEMO_CHECKLIST_20261103.md}: demonstration and reproduction checklist.
\end{itemize}
\begin{verbatim}
. scripts/use_d_runtime.ps1
.venv/Scripts/python.exe scripts/reaudit_recovered_validation.py
.venv/Scripts/python.exe -m pytest -q
\end{verbatim}
The group template, roster, FPT branding, orange section style, table of contents,
header/footer and bibliography are preserved. The compiler/cache and pre-edit
template backup are on D:. No external submission has been performed.
